\section{Тест производительности}
Для анализа эффективности реализованного алгоритма КМП было проведено сравнение с высокооптимизированным стандартным поиском \texttt{std::string::find()}. 
Тестирование проводилось на различных объемах данных: от 0.5 до 2 миллионов слов.
\begin{alltt}
    Тест 1: 500,000 слов
    tim@tim-swift:~/.../$ ./std < ./input.txt
    std::string::find() Time: 37.4529 ms
    tim@tim-swift:~/.../$ ./my < ./input.txt
    KMP Time: 57.269 ms
    
    Тест 2: 1,000,000 слов
    tim@tim-swift:~/.../$ ./std < ./input.txt
    std::string::find() Time: 47.0277 ms
    tim@tim-swift:~/.../$ ./my < ./input.txt
    KMP Time: 110.918 ms
    
    Тест 3: 2,000,000 слов
    tim@tim-swift:~/.../$ ./std < ./input.txt
    std::string::find() Time: 97.3903 ms
    tim@tim-swift:~/.../$ ./my < ./input.txt
    KMP Time: 202.918 ms
\end{alltt}

\textbf{Сводная таблица результатов:}

\begin{table}[h!]
    \centering
    % Фиксируем ширину первой колонки (5.5см), остальные авто-подгон по центру
    \begin{tabular}{|p{5.5cm}|c|c|c|}
        \hline
        Количество слов & 500 000 & 1 000 000 & 2 000 000 \\ \hline
        \texttt{std::string::find()} & 37.45 ms & 47.03 ms & 97.39 ms \\ \hline
        Алгоритм КМП (авторский) & 57.27 ms & 110.92 ms & 202.92 ms \\ \hline
    \end{tabular}
\end{table}
\textbf{Анализ результатов:}
Наблюдается четкая линейная масштабируемость обоих алгоритмов. При увеличении объема входных данных в 4 раза (с 0.5 млн до 2 млн слов), время выполнения увеличивается пропорционально, что подтверждает линейную сложность для обоих подходов.

\textbf{Преимущество стандартной библиотеки:} \texttt{std::string::find()} стабильно работает быстрее (в 1.5--2.3 раза). Это обусловлено использованием низкоуровневых оптимизаций:
\begin{itemize}
    \item \textbf{SIMD-инструкции:} Стандартная функция сравнивает сразу несколько символов за один такт процессора.
    \item \textbf{Отсутствие накладных расходов:} Моя реализация тратит время на пословную токенизацию, создание объектов \texttt{std::string} и работу с кольцевым буфером, в то время как \texttt{std::string::find()} выполняет поиск в "сыром" бинарном буфере.
\end{itemize}
\textbf{Особенности реализации:} Несмотря на отставание в скорости, алгоритм КМП производит поиск именно по логическим единицам (словам) и использует константный объем дополнительной памяти, что критически важно при работе с потоковыми данными большого объема, которые невозможно целиком загрузить в ОЗУ.